\documentclass[11pt]{article}
\usepackage[margin=1.1in]{geometry}
\usepackage{amsmath,amssymb,amsthm,booktabs,longtable,xcolor,array}
\usepackage[colorlinks=true,linkcolor=blue!50!black,citecolor=blue!50!black,urlcolor=blue!50!black]{hyperref}
\usepackage{orcidlink}
\usepackage{graphicx}
\graphicspath{{paper/figures/}}

\newtheorem{fact}{Fact}
\newtheorem{prop}{Proposition}
\newtheorem{conj}{Conjecture}
\newtheorem{thm}{Theorem}
\newtheorem{lem}{Lemma}
\newtheorem{cor}{Corollary}
\theoremstyle{definition}
\newtheorem{defn}{Definition}
\theoremstyle{remark}
\newtheorem{remark}{Remark}
\newcommand{\lam}{\lambda}
\newcommand{\code}[1]{\texttt{#1}}
\emergencystretch=2em
\hyphenation{re-flec-tion-ad-mit-ting quasi-crys-tal-line}

\title{A Flat Initial Surface? A Speculative Reading of the Spectre and Chair44 Tilings}

\author{Brent Hartshorn \orcidlink{0009-0004-2853-655X}
        and Tristan Romanov}

\begin{document}
\maketitle


\begin{abstract}
A conjecture about the initial state of the universe is stated so that it can be argued with: the Big Bang was a flat
two-dimensional singularity, connected to the planar Spectre monotile, and the emergence of three-dimensional space is
connected to Tsiokos's Chair44.  Two classical background facts, a reading through the inflation arithmetic of the two
tilings, and a proposal on the Weyl tensor are given.  Nothing is derived; what would support or retire the reading is
stated.
\end{abstract}


\section{The conjecture}
\label{sec:thesis}

This note states the hypothesis that motivated the companion paper~\cite{HartshornRomanov}, which compares the Spectre
monotile~\cite{SmithEtAl2} and Tsiokos's Chair44~\cite{Tsiokos1} by exact computation.  Nothing computed there depends
on it, and nothing computed there tests it.  It is stated here so that it can be argued with.

\begin{conj}[The initial surface]
\label{conj:bang}
The Big Bang was not a point-like singularity but a flat two-dimensional one.  That flat surface is connected to the
Spectre monotile, and the emergence of three-dimensional space from it is connected to the Chair44.
\end{conj}

\noindent
Two background facts are kept apart from the conjecture.  The first is classical and exact.  An infinite uniform flat sheet
of surface density $\sigma$ produces a gravitational field $g=2\pi G\sigma$ that does not depend on the distance from the
sheet, because the mass in view grows as $r^{2}$ exactly as the force per unit mass falls as $r^{-2}$ (equivalently, by
Gauss's law on a pillbox).  Fields of parallel sheets are uniform and superpose, so sheets of equal and opposite density
cancel outside, and two equal sheets cancel between them.  The second concerns the setting.  Ghost-free bimetric gravity
of Hassan and Rosen~\cite{HassanRosen} gives dynamics to a second metric, and a pre-geometric origin for it has been
proposed~\cite{PreGeometric}; a picture in which the initial state consists of two sheets, one for each metric, and the Big
Bang is their separation is the kind of reading that such a setting invites.  No such model is built here, and no claim
is made about the dual-sheet cosmologies themselves.  The earlier observation that the area law
$S_{\rm rel}=\delta A/4G$ is consistent with the $s$-wave dynamics of the Nariai horizon~\cite{hartshorn2026}, the converse
of the result of Dorau and Much~\cite{DorauMuch}, is the setting in which the question arose.

What the computations of~\cite{HartshornRomanov} offer is a precise form of the contrast the conjecture needs.  A planar tile with a unimodular
inflation, here the Spectre, has an internal space of dimension two and a window; a tile with an integer inflation, here the
Chair44, has none and is $2$-adic.  The conjecture would then say that the surface is of the first kind and the bulk of the
second, and that the passage from one to the other is the emergence of the third dimension.  That is a reading and not a
result: no map from tile data to a physical quantity is given, the dichotomy is not derived from any dynamics, and the
passage between the two kinds of inflation is not constructed.  The hierarchy section of~\cite{HartshornRomanov} tried the substitution hierarchy as that quantity, and it does not
supply the map.  The reading would
be supported by a physical quantity computed from the tile data that distinguishes the two cases, and retired by showing
that the dichotomy is an artefact of the two examples, for instance by a three-dimensional tile with a unimodular
inflation.  A planar counterexample of the other kind already exists: the hexagonal monotile of Socolar and
Taylor~\cite{SocolarTaylor} inflates by $2$, and its tilings are model sets with a $2$-adic internal
space~\cite{LeeMoodyTS}.  Its matching rules are carried by decorations, or by a shape that is not a connected region, so
the dichotomy can survive it only as a statement about connected tiles whose shape alone forces aperiodicity; this
narrowing is part of the reading, not a consequence of it.

\begin{remark}[A proposal on the Weyl tensor]
\label{rem:weyl}
In three spacetime dimensions the Weyl tensor vanishes identically, the Riemann tensor is fixed by the Ricci tensor, and
vacuum general relativity has no local gravitational degrees of freedom: gravitational waves need a third spatial
dimension.  That is classical.  The \emph{proposal}, the first author's and not a derived result, is that this is the form in which
a flat two-dimensional initial surface would have vanishing Weyl curvature (compare Penrose's Weyl curvature hypothesis, that
the Weyl tensor vanishes at the initial singularity): a $2{+}1$ substrate has no Weyl tensor to set to zero, and the Weyl
degrees of freedom would first appear with the third spatial dimension, the step the conjecture assigns to the Chair44.
Nothing here derives it.  No map from tile data to a metric is given; it is not shown that a $2{+}1$ substrate can acquire
$3{+}1$ Weyl degrees of freedom as space emerges; and the hierarchy section of~\cite{HartshornRomanov} leaves open what
the hierarchy depth is: if it is the
third spatial dimension, a separate time is needed, and if it is look-back time, gravitational waves need a propagation rule that
the tiling does not supply.
\end{remark}


\begin{thebibliography}{9}

\bibitem{HartshornRomanov} B.~Hartshorn and T.~Romanov, \emph{The Internal Space and Boundary Growth of the Spectre and Tsiokos's Chair44 Tilings}, preprint (2026),
\url{https://github.com/brentharts/spectre}.

\bibitem{SmithEtAl2} D.~Smith, J.~S.~Myers, C.~S.~Kaplan and C.~Goodman-Strauss, \emph{A chiral aperiodic monotile},
Combin.\ Theory \textbf{4}(2) (2024); arXiv:2305.17743.

\bibitem{Tsiokos1} I.~Tsiokos, \emph{A strongly aperiodic monotile in three dimensions} (2026),
arXiv:2609.19214.

\bibitem{SocolarTaylor} J.~E.~S.~Socolar and J.~M.~Taylor, \emph{An aperiodic hexagonal tile}, J.\ Combin.\ Theory Ser.\ A
\textbf{118} (2011) 2207--2231; arXiv:1003.4279.

\bibitem{LeeMoodyTS} J.-Y.~Lee and R.~V.~Moody, \emph{Taylor--Socolar hexagonal tilings as model sets}, Symmetry
\textbf{5} (2013) 1--46; arXiv:1207.6237.


\bibitem{HassanRosen} S.~F.~Hassan and R.~A.~Rosen, \emph{Bimetric gravity from ghost-free massive gravity},
J.\ High Energy Phys.\ \textbf{2012} (2012) 126; \url{https://link.springer.com/article/10.1007/JHEP02(2012)126}

\bibitem{PreGeometric} A.~Addazi et al., \emph{Pre-geometric origin of ghost-free massive bigravity}, Phys.\ Lett.\ B
\textbf{882} (2026) 140989; \url{https://doi.org/10.1016/j.physletb.2026.140989}

\bibitem{hartshorn2026} B.~Hartshorn,
\emph{Gravity from relative entropy: Jackiw--Teitelboim dynamics on the Nariai horizon}, preprint (2026),
\url{https://doi.org/10.21203/rs.3.rs-10426867/v1}

\bibitem{DorauMuch} P.~Dorau and A.~Much, \emph{From quantum relative entropy to the semiclassical Einstein equations}, Phys.\ Rev.\ Lett.\ \textbf{136} (2026) 091602; arXiv:2510.24491;
\url{https://doi.org/10.1103/lmq8-nsty}



\end{thebibliography}

\end{document}
